\section{Проектирование и разработка программного средства}

Третий раздел посвящен проектированию и разработке программного средства сквозной аналитики для \eng{e-commerce} проектов. На этом этапе требования, сформулированные во втором разделе, переводятся в конкретные архитектурные, технологические и интерфейсные решения. Необходимо описать, каким образом клиентская и серверная части взаимодействуют между собой, как организована модель данных, какие экранные формы получает пользователь и какие механизмы обеспечивают корректность обработки аналитической информации.

Для дипломного проекта особенно важно сохранить баланс между инженерной целостностью и воспроизводимостью. Система должна быть достаточно реалистичной, чтобы продемонстрировать полный цикл работы с аналитикой \eng{e-commerce} проекта, но при этом не должна зависеть от внешнего интернет-магазина, сторонних учетных записей или коммерческих интеграций. Поэтому проектирование ведется вокруг автономной \eng{web}-системы, где данные вводятся через интерфейс или формируются в виде демонстрационного массива.

В рамках раздела последовательно рассматриваются архитектурные решения и выбор технологий, проектирование пользовательского интерфейса, разработка модели данных, организация бизнес-логики и спецификация \eng{REST API} с механизмами информационной безопасности. Такой порядок соответствует переходу от общего устройства системы к ее конкретным компонентам и внутренним ограничениям.

\subsection{Архитектурные решения и технологии реализации программного средства}

Разрабатываемое программное средство представляет собой распределенную \eng{web}-систему, в которой пользовательский интерфейс, серверная логика и база данных разделены на самостоятельные уровни. Такое разделение позволяет независимо развивать клиентскую и серверную части, а также упрощает сопровождение и проверку отдельных компонентов. Архитектура программного средства сквозной аналитики представлена на рисунке~\ref{fig:analytics-architecture}.

\begin{figure}[H]
    \centering
    \diplomaimage{images/fig-3-1-analytics-architecture.png}{100mm}
    \caption{Архитектура программного средства сквозной аналитики}
    \label{fig:analytics-architecture}
\end{figure}

На стороне клиента используется приложение на базе \eng{Next.js 15}, \eng{React 19} и \eng{TypeScript}. Клиентская часть отвечает за отображение форм заявок, табличного ввода, \eng{dashboard}, отчетов, рекомендаций и административных страниц. На стороне сервера применяется \eng{Go 1.23}, маршрутизатор \eng{chi}, драйвер \eng{pgx} и генератор \eng{sqlc}. Серверная часть реализует аутентификацию, управление ролями, работу с проектами и расчет аналитических показателей. В качестве основной реляционной базы данных используется \eng{PostgreSQL 16}. Развертывание приложения выполняется через \eng{Docker Compose}, а публикация \eng{web}-системы и завершение \eng{HTTPS}-соединений реализуются через \eng{Caddy}.

Такое построение системы соответствует требованиям предметной области. Клиенту и \eng{SEO}-специалисту необходим быстрый доступ к аналитическим экранам и табличным формам, а серверной части -- четко структурированная прикладная логика и надежная работа с транзакционными данными. Выбор \eng{Go} для серверной части обусловлен высокой производительностью, явной моделью обработки запросов и компактностью исполняемой среды. Выбор \eng{Next.js} и \eng{React} для клиентской части обусловлен компонентной архитектурой, удобством построения сложных интерфейсов и поддержкой современных сценариев \eng{SPA}/\eng{SSR}.

До окончательного выбора стека был выполнен сравнительный анализ языков программирования для серверной части. Его результаты приведены в таблице~\ref{tab:backend-languages}.

\begin{table}[H]
    \caption{Сравнительный анализ языков программирования для серверной части}
    \label{tab:backend-languages}
    {\small
    \setlength{\tabcolsep}{4pt}
    \begin{tabularx}{\textwidth}{|P{0.18\textwidth}|L|L|L|}
        \hline
        Язык & Преимущества & Ограничения & Вывод для проекта \\
        \hline
        \eng{Go 1.23} & Высокая производительность, простая модель конкурентности, компактный исполняемый сервис, удобная разработка \eng{REST API} & Менее богатая экосистема готовых \eng{ORM} по сравнению с динамическими языками & Выбран как основной язык серверной части \\
        \hline
        \eng{Rust} & Высокая производительность, строгая безопасность памяти, современная экосистема серверной разработки & Более высокий порог входа и большая трудоемкость реализации прикладной логики для учебного проекта & Перспективен, но избыточен по сложности для выбранной задачи \\
        \hline
        \eng{Node.js} & Единый язык с клиентской частью, развитая экосистема пакетов, удобство работы с \eng{JSON} & Более высокая сложность контроля архитектуры при росте проекта, зависимость от большого числа пакетов & Может использоваться для \eng{web}-сервисов, но уступает выбранному стеку по строгости серверной реализации \\
        \hline
        \eng{Java} & Зрелая корпоративная экосистема, типизация, надежные инструменты разработки & Более тяжелая инфраструктура для учебного проекта, высокий порог детализации конфигурации & Избыточен для поставленной задачи \\
        \hline
    \end{tabularx}
    }
\end{table}

Для организации серверных маршрутов и \eng{middleware} был выполнен сравнительный анализ фреймворков и библиотек для разработки серверной части. Результаты приведены в таблице~\ref{tab:backend-frameworks}.

\begin{table}[H]
    \caption{Сравнительный анализ фреймворков для разработки}
    \label{tab:backend-frameworks}
    {\small
    \setlength{\tabcolsep}{4pt}
    \begin{tabularx}{\textwidth}{|P{0.18\textwidth}|L|L|L|}
        \hline
        Решение & Преимущества & Ограничения & Вывод для проекта \\
        \hline
        \eng{chi} & Легковесная маршрутизация, прозрачная работа с \eng{middleware}, хорошая совместимость со стандартной библиотекой \eng{Go} & Требует явной организации слоев приложения & Выбран для серверной части \\
        \hline
        \eng{Gin} & Высокая популярность, большое число готовых расширений, быстрый старт & Более выраженная зависимость от собственного стиля фреймворка & Подходит, но менее нейтрален по архитектуре \\
        \hline
        \eng{Echo} & Удобный набор компонентов для \eng{REST}-разработки, встроенные возможности валидации и \eng{middleware} & Менее предпочтителен при желании сохранять близость к стандартной библиотеке & Может применяться, но не выбран \\
        \hline
        \eng{Fiber} & Высокая производительность и лаконичный синтаксис & Отличается от стандартного \eng{net/http}, что повышает связанность проекта с конкретным стеком & Нецелесообразен для учебного проекта \\
        \hline
    \end{tabularx}
    }
\end{table}

Архитектурная организация системы на уровне внутренних слоев соответствует паттерну \eng{MVC}, адаптированному к распределенной \eng{web}-системе. В клиентской части слой представления формируется \eng{React}-компонентами, в серверной части контроллеры представлены обработчиками \eng{HTTP}-запросов, а слой модели включает сущности предметной области и прикладную логику. Архитектурный паттерн \eng{MVC} представлен на рисунке~\ref{fig:mvc-pattern}.

\begin{figure}[H]
    \centering
    \diplomaimage{images/fig-3-2-mvc-pattern.png}{96mm}
    \caption{Архитектурный паттерн \eng{MVC}}
    \label{fig:mvc-pattern}
\end{figure}

Для клиентской части также был выполнен сравнительный анализ технологий интерфейса. Результаты приведены в таблице~\ref{tab:frontend-technologies}.

\begin{table}[H]
    \caption{Сравнительный анализ технологий клиентской части}
    \label{tab:frontend-technologies}
    {\small
    \setlength{\tabcolsep}{4pt}
    \begin{tabularx}{\textwidth}{|P{0.22\textwidth}|L|L|L|}
        \hline
        Технология & Преимущества & Ограничения & Вывод для проекта \\
        \hline
        \eng{Next.js 15 + React 19 + TypeScript} & Компонентная архитектура, маршрутизация, удобное построение сложных интерфейсов, строгая типизация & Требует дисциплины при проектировании состояния и взаимодействия с \eng{API} & Выбран как основной стек клиентской части \\
        \hline
        \eng{Vue} & Удобный порог входа, реактивность, компактный синтаксис & В проекте нет необходимости менять базовый стек при уже выбранной экосистеме \eng{React} & Подходит, но не выбран \\
        \hline
        \eng{Angular} & Целостный фреймворк, развитая типизация, строгая структура приложений & Более тяжеловесен для учебной \eng{SaaS}-системы с ограниченным объемом разработки & Избыточен для проекта \\
        \hline
        Классический \eng{HTML/CSS/JS} & Простота развертывания и минимальный стек & Недостаточно удобен для масштабируемого интерфейса с большим числом экранов и состояний & Не подходит для целевой системы \\
        \hline
    \end{tabularx}
    }
\end{table}

Для описания границ системы и взаимодействия с внешними участниками используется модель \eng{C4}. На контекстном уровне она показывает, какие пользователи взаимодействуют с сервисом сквозной аналитики, и какие внешние инфраструктурные элементы участвуют в обработке данных. Контекстная диаграмма \eng{C4} системы аналитики представлена на рисунке~\ref{fig:c4-context}.

\begin{figure}[H]
    \centering
    \diplomaimage{images/fig-3-3-c4-context.png}{100mm}
    \caption{Контекстная диаграмма \eng{C4} системы аналитики}
    \label{fig:c4-context}
\end{figure}

На контейнерном уровне система разбивается на \eng{web}-клиент, серверный \eng{API}, сервис аналитики, сервисы управления проектами и пользователями, а также базу данных. Такое деление позволяет явным образом показать распределение ответственности между компонентами и подготавливает основу для описания прикладной логики. Диаграмма контейнеров системы представлена на рисунке~\ref{fig:c4-containers}.

\begin{figure}[H]
    \centering
    \diplomaimage{images/fig-3-4-c4-containers.png}{104mm}
    \caption{Диаграмма контейнеров системы}
    \label{fig:c4-containers}
\end{figure}

Сравнение решений развертывания и публикации интерфейса приведено в таблице~\ref{tab:deployment-technologies}.

\begin{table}[H]
    \caption{Сравнительный анализ технологий развертывания и публикации}
    \label{tab:deployment-technologies}
    {\small
    \setlength{\tabcolsep}{4pt}
    \begin{tabularx}{\textwidth}{|P{0.20\textwidth}|L|L|L|}
        \hline
        Технология & Преимущества & Ограничения & Вывод для проекта \\
        \hline
        \eng{Docker Compose + Caddy} & Воспроизводимое развертывание, простой \eng{reverse proxy}, удобная работа с \eng{HTTPS} и статическими ресурсами & Требует базовой контейнерной инфраструктуры & Выбран для демонстрационного стенда \\
        \hline
        \eng{Docker Compose + Traefik} & Удобная контейнерная интеграция, автоматизация маршрутизации, поддержка \eng{TLS} & Избыточен для локального учебного стенда и сложнее в настройке, чем выбранное решение & Возможен, но не выбран \\
        \hline
        Локальный запуск без контейнеров & Быстрый старт на одной машине & Низкая воспроизводимость среды и зависимость от локальной конфигурации & Не соответствует требованиям демонстрации \\
        \hline
        Облачная публикация \eng{PaaS} & Удобное внешнее размещение и автоматизация развертывания & Избыточна для учебного проекта, добавляет внешние зависимости & Не требуется в рамках диплома \\
        \hline
    \end{tabularx}
    }
\end{table}

Наиболее детальное представление серверной части дается на уровне компонентной диаграммы. Здесь выделяются контроллеры, сервисы прикладной логики, репозитории и сущности. Такой подход позволяет показать, что система не сводится к набору \eng{HTTP}-обработчиков, а имеет внутреннее разделение ответственности. Диаграмма компонентов \eng{backend} представлена на рисунке~\ref{fig:backend-components}.

\begin{figure}[H]
    \centering
    \diplomaimage{images/fig-3-5-backend-components.png}{104mm}
    \caption{Диаграмма компонентов \eng{backend}}
    \label{fig:backend-components}
\end{figure}

Таким образом, архитектурные решения раздела 3.1 показывают, что выбранный стек и структура приложения соответствуют поставленным требованиям. Клиентская часть обеспечивает удобную работу с интерфейсом, серверная часть реализует прикладную и аналитическую логику, а база данных хранит связанный набор сущностей предметной области. Выбранные технологии ориентированы на воспроизводимость, типизацию, явную структуру слоев и удобство демонстрации дипломного проекта.

\subsection{Проектирование и разработка пользовательского интерфейса}

Пользовательский интерфейс программного средства разрабатывается как рабочая \eng{B2B SaaS}-система для аналитического сопровождения \eng{e-commerce} проектов. Интерфейс не является маркетинговым сайтом: он ориентирован на ежедневную работу с заявками, проектами, таблицами, графиками, отчетами и административными функциями. Основной упор делается на понятность структуры, плотность полезной информации и предсказуемость пользовательских действий.

Визуальная концепция строится на светлой теме, нейтральном фоне, синем акценте и аккуратной сетке рабочих блоков. Такой стиль соответствует задачам аналитического продукта, где пользователь должен быстро считывать показатели, статусы и изменения данных. Базовые элементы оформления: карточки \eng{KPI}, табличные блоки, графики, фильтры, формы ввода, статусные бейджи и навигационные панели.

Согласованность визуальных элементов фиксируется в \eng{design system}, которая объединяет набор цветов, типографику, кнопки, поля ввода, графические элементы и правила отображения статусов. \eng{Design System} интерфейса аналитической платформы представлен на рисунке~\ref{fig:design-system}.

\begin{figure}[H]
    \centering
    \diplomaimage{images/fig-3-6-design-system.png}{92mm}
    \caption{\eng{Design System} интерфейса аналитической платформы}
    \label{fig:design-system}
\end{figure}

Для проектирования интерфейса важно учитывать различие пользовательских ролей. Клиенту необходимы сценарии подачи заявки, заполнения брифа, загрузки табличных данных и просмотра аналитического результата. \eng{SEO}-специалисту требуются рабочие экраны проверки данных, анализа источников и подготовки отчета. Администратору нужны страницы управления пользователями, проектами, резервным копированием и журналом действий. Поэтому интерфейс проектируется как единая система с разными рабочими контурами.

Пользовательский сценарий клиента начинается с подачи заявки и заканчивается просмотром отчета. На промежуточных этапах клиент передает сведения о бизнесе, заполняет таблицы событий и заказов, а затем анализирует рассчитанные показатели по проекту. \eng{User Flow} клиента представлен на рисунке~\ref{fig:user-flow-client}.

\begin{figure}[H]
    \centering
    \diplomaimage{images/fig-3-7-client-user-flow.png}{96mm}
    \caption{\eng{User Flow} клиента}
    \label{fig:user-flow-client}
\end{figure}

Пользовательский сценарий \eng{SEO}-специалиста отличается большей насыщенностью прикладными операциями. Специалист принимает заявку, проводит диагностику проекта, настраивает правила аналитики, проверяет качество данных, интерпретирует показатели и публикует отчет. \eng{User Flow} \eng{SEO}-специалиста представлен на рисунке~\ref{fig:user-flow-seo}.

\begin{figure}[H]
    \centering
    \diplomaimage{images/fig-3-8-seo-user-flow.png}{96mm}
    \caption{\eng{User Flow} \eng{SEO}-специалиста}
    \label{fig:user-flow-seo}
\end{figure}

Навигационная структура интерфейса основана на левом боковом меню и верхней панели выбора проекта и периода. Такое решение обеспечивает быстрый переход между связанными разделами и удобно для работы с несколькими проектами. Навигационное меню системы аналитики представлено на рисунке~\ref{fig:navigation-menu}.

\begin{figure}[H]
    \centering
    \diplomaimage{images/fig-3-9-navigation-menu.png}{68mm}
    \caption{Навигационное меню системы аналитики}
    \label{fig:navigation-menu}
\end{figure}

Центральным экраном пользовательского интерфейса является главная \eng{dashboard}-страница проекта. Она отображает ключевые показатели, динамику трафика, конверсии и выручки, а также последние аналитические выводы. Такой экран должен позволять специалисту быстро оценить текущее состояние проекта, а клиенту -- получить целостную картину результата. Главная \eng{dashboard}-страница представлена на рисунке~\ref{fig:main-dashboard}.

\begin{figure}[H]
    \centering
    \diplomaimage{images/fig-3-10-main-dashboard.png}{105mm}
    \caption{Главная \eng{dashboard}-страница}
    \label{fig:main-dashboard}
\end{figure}

Страница проектов используется для управления активными \eng{e-commerce} проектами. На ней отображаются основные сведения о сайте, клиенте, назначенном специалисте, статусе, подключенных источниках и ключевых показателях. Этот экран важен для систематизации нескольких проектов в едином интерфейсе. Страница проектов \eng{e-commerce} представлена на рисунке~\ref{fig:ecommerce-projects}.

\begin{figure}[H]
    \centering
    \diplomaimage{images/fig-3-11-ecommerce-projects.png}{105mm}
    \caption{Страница проектов \eng{e-commerce}}
    \label{fig:ecommerce-projects}
\end{figure}

Страница аналитических задач отражает рабочее планирование \eng{SEO}-специалиста: диагностику, проверку данных, пересчет показателей, подготовку отчетов и дедлайны по проектам. Наличие такого экрана демонстрирует, что система поддерживает не только хранение данных, но и организацию аналитической работы. Страница аналитических задач представлена на рисунке~\ref{fig:analytics-tasks}.

\begin{figure}[H]
    \centering
    \diplomaimage{images/fig-3-12-analytics-tasks.png}{105mm}
    \caption{Страница аналитических задач}
    \label{fig:analytics-tasks}
\end{figure}

Страница клиентов выполняет функции упрощенного \eng{CRM}-раздела внутри аналитической системы. Она объединяет сведения о клиентах, связанных проектах, статусах работы и аналитической активности. Страница клиентов представлена на рисунке~\ref{fig:clients-page}.

\begin{figure}[H]
    \centering
    \diplomaimage{images/fig-3-13-clients.png}{105mm}
    \caption{Страница клиентов}
    \label{fig:clients-page}
\end{figure}

Для предметной области сквозной аналитики особенно важен специализированный экран источников трафика. На нем пользователь видит показатели по \eng{organic}, \eng{paid\_search}, \eng{social}, \eng{direct}, \eng{referral} и \eng{email}, а также может сравнить каналы по заказам, выручке, конверсии, расходам и \eng{ROMI}. Страница источников трафика представлена на рисунке~\ref{fig:traffic-sources-page}.

\begin{figure}[H]
    \centering
    \diplomaimage{images/fig-3-14-traffic-sources.png}{105mm}
    \caption{Страница источников трафика}
    \label{fig:traffic-sources-page}
\end{figure}

Финальным пользовательским результатом выступает экран аналитических отчетов. Он объединяет сведения о периоде анализа, рассчитанных показателях, рекомендациях и статусе публикации. За счет этого отчет становится не отдельным файлом, а элементом управляемого процесса внутри системы. Страница аналитических отчетов представлена на рисунке~\ref{fig:analytics-reports-page}.

\begin{figure}[H]
    \centering
    \diplomaimage{images/fig-3-15-analytics-reports.png}{105mm}
    \caption{Страница аналитических отчетов}
    \label{fig:analytics-reports-page}
\end{figure}

Таким образом, пользовательский интерфейс системы проектируется как рабочее пространство для трех ролей и отражает прикладную логику сквозной аналитики. Он объединяет процессы регистрации заявок, управления проектами, загрузки данных, анализа \eng{KPI}, сравнения источников и публикации отчетов. Такая структура интерфейса соответствует задачам дипломного проекта и обеспечивает достаточную детализацию для демонстрации программного средства.

\subsection{Разработка модели данных}

Модель данных программного средства должна отражать как организационные аспекты сопровождения аналитики, так и собственно аналитические сущности. Для этого в базе данных выделяются пользователи, заявки, проекты, источники трафика, события, заказы, возвраты, отчеты, рекомендации, метрики, резервные копии и журнал действий. Связи между сущностями обеспечивают возможность перехода от заявки клиента к итоговому аналитическому отчету в рамках единой реляционной модели.

На концептуальном уровне основой модели выступает сущность проекта. С ней связываются источники трафика, события, заказы, аналитические показатели, отчеты и рекомендации. Заявка может выступать исходной точкой создания проекта, а пользовательские роли определяют права на просмотр и изменение этих данных. \eng{ER}-диаграмма базы данных системы представлена на рисунке~\ref{fig:er-database}.

\begin{figure}[H]
    \centering
    \diplomaimage{images/fig-3-16-er-database.png}{105mm}
    \caption{\eng{ER}-диаграмма базы данных системы}
    \label{fig:er-database}
\end{figure}

При проектировании модели данных учитывались требования к нормализации и повторному использованию аналитических сведений. Выделение отдельных таблиц для проектов, источников, событий, заказов и возвратов позволяет хранить данные без избыточного дублирования и рассчитывать показатели в разных разрезах. Дополнительные служебные таблицы, такие как журнал действий и резервные копии, обеспечивают контроль и эксплуатацию системы.

Описание сущностей и атрибутов приведено в таблице~\ref{tab:entities-attributes}.

{\small
\setlength{\tabcolsep}{4pt}
\begin{longtable}{|P{0.19\textwidth}|P{0.33\textwidth}|P{0.38\textwidth}|}
    \caption{Описание сущностей и атрибутов}
    \label{tab:entities-attributes}\\
    \hline
    Сущность & Основные атрибуты & Назначение \\
    \hline
    \endfirsthead
    \caption*{Продолжение таблицы \thetable}\\
    \hline
    Сущность & Основные атрибуты & Назначение \\
    \hline
    \endhead
    \eng{users} & \eng{id}, \eng{full\_name}, \eng{email}, \eng{password\_hash}, \eng{role}, \eng{is\_active}, \eng{created\_at} & Хранение учетных записей клиентов, \eng{SEO}-специалистов и администраторов \\
    \hline
    \eng{applications} & \eng{id}, \eng{client\_id}, \eng{project\_id}, \eng{seo\_specialist\_id}, \eng{description}, \eng{goal}, \eng{status}, \eng{created\_at} & Регистрация заявок на подключение аналитики \\
    \hline
    \eng{projects} & \eng{id}, \eng{client\_id}, \eng{seo\_specialist\_id}, \eng{name}, \eng{site\_url}, \eng{niche}, \eng{status}, \eng{created\_at} & Основная сущность аналитического проекта \\
    \hline
    \eng{traffic\_sources} & \eng{id}, \eng{project\_id}, \eng{name}, \eng{source\_type}, \eng{is\_active} & Хранение каналов и источников трафика проекта \\
    \hline
    \eng{events} & \eng{id}, \eng{project\_id}, \eng{traffic\_source\_id}, \eng{event\_name}, \eng{occurred\_at}, \eng{anonymous\_id} & Фиксация пользовательских событий на сайте \\
    \hline
    \eng{orders} & \eng{id}, \eng{project\_id}, \eng{traffic\_source\_id}, \eng{amount}, \eng{status}, \eng{ordered\_at} & Хранение заказов и сумм покупок \\
    \hline
    \eng{refunds} & \eng{id}, \eng{order\_id}, \eng{amount}, \eng{refunded\_at}, \eng{reason} & Фиксация возвратов и их причин \\
    \hline
    \eng{reports} & \eng{id}, \eng{project\_id}, \eng{author\_id}, \eng{period\_start}, \eng{period\_end}, \eng{status}, \eng{created\_at} & Хранение аналитических отчетов по проекту \\
    \hline
    \eng{recommendations} & \eng{id}, \eng{project\_id}, \eng{report\_id}, \eng{title}, \eng{description}, \eng{priority}, \eng{expected\_effect}, \eng{status}, \eng{created\_at} & Хранение рекомендаций, сформированных по результатам анализа \\
    \hline
    \eng{analytics\_metrics} & \eng{id}, \eng{project\_id}, \eng{period\_start}, \eng{period\_end}, \eng{visits\_count}, \eng{orders\_count}, \eng{revenue}, \eng{refunds\_amount}, \eng{conversion\_rate}, \eng{average\_order\_value} & Хранение рассчитанных агрегированных показателей \\
    \hline
    \eng{audit\_logs} & \eng{id}, \eng{user\_id}, \eng{action}, \eng{entity\_type}, \eng{entity\_id}, \eng{created\_at}, \eng{ip\_address} & Журналирование действий пользователей \\
    \hline
    \eng{backups} & \eng{id}, \eng{file\_name}, \eng{status}, \eng{created\_at}, \eng{created\_by} & Учет операций резервного копирования \\
    \hline
\end{longtable}
}

Приведенная модель данных обеспечивает возможность последовательного прохождения жизненного цикла проекта: от входящей заявки до опубликованного отчета. Кроме того, она поддерживает предметные требования по расчету переходов, заказов, выручки, возвратов, конверсии и показателей эффективности источников трафика. За счет выделения отдельных сущностей модель остается расширяемой и пригодной для дальнейшего развития.

\subsection{Описание организации бизнес-логики программного средства}

Бизнес-логика программного средства связывает пользовательские действия, модель данных и внутренние вычислительные процедуры. Она определяет, как заявка превращается в проект, каким образом проверяются исходные данные, когда выполняется перерасчет аналитических показателей и как формируется итоговый отчет. Именно этот слой превращает систему из набора экранов и таблиц в связанное прикладное решение.

Первый ключевой процесс связан с обработкой аналитической заявки. После подачи заявки клиентом система сохраняет ее, назначает ответственного \eng{SEO}-специалиста и позволяет перевести заявку в рабочий статус. Далее специалист проводит диагностику, формирует проект и инициирует этап сбора данных. Диаграмма деятельности процесса обработки аналитической заявки представлена на рисунке~\ref{fig:application-processing-activity}.

\begin{figure}[H]
    \centering
    \diplomaimage{images/fig-3-17-application-processing-activity.png}{102mm}
    \caption{Диаграмма деятельности процесса обработки аналитической заявки}
    \label{fig:application-processing-activity}
\end{figure}

После создания проекта начинается цикл аналитической обработки данных. Клиент или администратор передают данные о переходах, событиях, заказах и возвратах. Серверная часть выполняет валидацию структуры и обязательных полей, затем аналитический модуль рассчитывает агрегированные показатели, а результаты становятся доступными для \eng{dashboard} и отчетов. Последовательность взаимодействия компонентов при расчете аналитики представлена на рисунке~\ref{fig:analytics-calculation-sequence}.

\begin{figure}[H]
    \centering
    \diplomaimage{images/fig-3-18-analytics-calculation-sequence.png}{106mm}
    \caption{Диаграмма последовательности расчета аналитики}
    \label{fig:analytics-calculation-sequence}
\end{figure}

Внутри прикладной логики можно выделить несколько устойчивых блоков. Блок управления заявками отвечает за регистрацию обращения, смену статусов и связь с проектом. Блок работы с данными проекта отвечает за ввод и проверку исходных сведений. Аналитический блок выполняет расчет \eng{KPI}, воронки продаж, эффективности источников и формирование итоговых метрик. Блок отчетности управляет публикацией отчета и фиксацией рекомендаций. Административный блок отвечает за пользователей, справочники, резервное копирование и демонстрационные данные.

Завершающий бизнес-процесс связан с подготовкой и публикацией отчета. На этом этапе \eng{SEO}-специалист анализирует рассчитанные показатели, формирует рекомендации по улучшению источников трафика и этапов воронки, после чего публикует отчет. Клиент получает доступ к результату и при необходимости может отправить отчет на доработку. Диаграмма деятельности процесса публикации отчета представлена на рисунке~\ref{fig:report-publication-activity}.

\begin{figure}[H]
    \centering
    \diplomaimage{images/fig-3-19-report-publication-activity.png}{102mm}
    \caption{Диаграмма деятельности процесса публикации отчета}
    \label{fig:report-publication-activity}
\end{figure}

Организация бизнес-логики показывает, что система поддерживает не только хранение данных, но и целостный процесс аналитического сопровождения \eng{e-commerce} проекта. Последовательность действий, статусов, проверок и расчетов обеспечивает воспроизводимый сценарий демонстрации и соответствует требованиям, сформулированным во втором разделе.

\subsection{Спецификация к разработанному API и механизмы обеспечения информационной безопасности}

Взаимодействие клиентской и серверной частей реализуется через \eng{REST API}, использующий формат \eng{JSON}. Сервер предоставляет точки доступа для авторизации, управления пользователями, заявками, проектами, табличными данными, аналитическими метриками, рекомендациями, отчетами, резервным копированием и журналом действий. Для документирования интерфейса используется \eng{OpenAPI}-описание, позволяющее формализовать структуру запросов и ответов без привязки к ручному описанию каждого маршрута.

Спецификация \eng{API} строится вокруг сущностей предметной области. Для заявок поддерживаются операции создания, просмотра, изменения статуса и назначения специалиста. Для проектов доступны операции получения карточки проекта, \eng{dashboard}, источников трафика и аналитических данных. Для ввода информации предусмотрены отдельные маршруты массовой загрузки событий, заказов и возвратов. Для отчетов реализуются операции подготовки, публикации, принятия и отправки на доработку. Такой набор маршрутов соответствует жизненному циклу проекта и позволяет клиентской части работать с системой без прямого доступа к базе данных.

С точки зрения информационной безопасности система должна обеспечивать аутентификацию пользователей, защиту паролей, разграничение прав и журналирование значимых действий. Аутентификация выполняется по \eng{email} и паролю, а доступ к защищенным маршрутам предоставляется по \eng{JWT}-токену. Пароли сохраняются только в виде криптографического хеша. Административные маршруты должны быть недоступны клиенту и \eng{SEO}-специалисту без соответствующей роли. Для внешнего доступа используется \eng{HTTPS}, завершаемый на уровне \eng{Caddy}.

Разграничение прав доступа в системе приведено в таблице~\ref{tab:access-rights}.

{\small
\setlength{\tabcolsep}{4pt}
\begin{longtable}{|P{0.19\textwidth}|P{0.31\textwidth}|P{0.40\textwidth}|}
    \caption{Разграничение прав доступа в системе}
    \label{tab:access-rights}\\
    \hline
    Роль & Разделы доступа & Основные разрешенные действия \\
    \hline
    \endfirsthead
    \caption*{Продолжение таблицы \thetable}\\
    \hline
    Роль & Разделы доступа & Основные разрешенные действия \\
    \hline
    \endhead
    Клиент & Заявки, бриф проекта, ввод данных, \eng{dashboard}, отчеты, рекомендации & Создание заявки, выбор \eng{SEO}-специалиста, заполнение данных о бизнесе, ввод событий и покупок, просмотр только собственных проектов и отчетов \\
    \hline
    \eng{SEO}-специалист & Назначенные заявки, проекты, аналитика, задачи, отчеты, рекомендации & Принятие заявки, диагностика проекта, настройка аналитики, проверка данных, анализ источников, публикация отчета \\
    \hline
    Администратор & Пользователи, роли, проекты, справочники, резервные копии, журнал действий, генерация данных & Управление пользователями и ролями, создание проектов, генерация демонстрационных данных, просмотр журнала, выполнение резервного копирования \\
    \hline
\end{longtable}
}

Дополнительными мерами защиты выступают проверка входных данных, контроль форматов и ограничение доступа к административным операциям. Для демонстрационного проекта также важно журналировать ключевые действия, чтобы можно было отследить публикацию отчета, изменение ролей, генерацию тестовых данных и создание резервных копий. Наличие таких механизмов делает систему не только функционально полной, но и инженерно состоятельной с точки зрения эксплуатации.

Таким образом, раздел 3 показывает, каким образом требования предметной области преобразуются в конкретную реализацию. Архитектура определяет состав компонентов и взаимодействий, интерфейс обеспечивает рабочие сценарии пользователей, модель данных хранит предметные сущности, бизнес-логика организует процесс сквозной аналитики, а \eng{REST API} и механизмы безопасности обеспечивают надежное функционирование системы.
